
\documentclass{article}
\usepackage{colortbl}
\usepackage{makecell}
\usepackage{multirow}
\usepackage{supertabular}

\begin{document}

\newcounter{utterance}

\centering \large Interaction Transcript for game `hot\_air\_balloon', experiment `air\_balloon\_survival\_de\_negotiation\_hard', episode 4 with gpt{-}oss{-}120b{-}t1.0.
\vspace{24pt}

{ \footnotesize  \setcounter{utterance}{1}
\setlength{\tabcolsep}{0pt}
\begin{supertabular}{c@{$\;$}|p{.15\linewidth}@{}p{.15\linewidth}p{.15\linewidth}p{.15\linewidth}p{.15\linewidth}p{.15\linewidth}}
    \# & $\;$A & \multicolumn{4}{c}{Game Master} & $\;\:$B\\
    \hline

    \theutterance \stepcounter{utterance}  
    & & \multicolumn{4}{p{0.6\linewidth}}{
        \cellcolor[rgb]{0.9,0.9,0.9}{
            \makecell[{{p{\linewidth}}}]{
                \texttt{\tiny{[P1$\langle$GM]}}
                \texttt{Du nimmst an einem kollaborativen Verhandlungspiel Teil.} \\
\\ 
\texttt{Zusammen mit einem anderen Teilnehmer musst du dich auf eine Reihe von Gegenständen entscheiden, die behalten werden. Jeder von euch hat eine persönliche Verteilung über die Wichtigkeit der einzelnen Gegenstände. Jeder von euch hat eine eigene Meinung darüber, wie wichtig jeder einzelne Gegenstand ist (Gegenstandswichtigkeit). Du kennst die Wichtigkeitsverteilung des anderen Spielers nicht. Zusätzlich siehst du, wie viel Aufwand jeder Gegenstand verursacht.  } \\
\texttt{Ihr dürft euch nur auf eine Reihe von Gegenständen einigen, wenn der Gesamtaufwand der ausgewählten Gegenstände den Maximalaufwand nicht überschreitet:} \\
\\ 
\texttt{Maximalaufwand = 5138} \\
\\ 
\texttt{Hier sind die einzelnen Aufwände der Gegenstände:} \\
\\ 
\texttt{Aufwand der Gegenstände = \{"C33": 903, "A89": 757, "A48": 850, "A04": 746, "A41": 524, "C64": 204, "C49": 946, "A61": 473, "C04": 616, "B53": 855, "B85": 530, "B61": 419, "A45": 960, "A03": 763, "B21": 730\}} \\
\\ 
\texttt{Hier ist deine persönliche Verteilung der Wichtigkeit der einzelnen Gegenstände:} \\
\\ 
\texttt{Werte der Gegenstandswichtigkeit = \{"C33": 138, "A89": 583, "A48": 868, "A04": 822, "A41": 783, "C64": 65, "C49": 262, "A61": 121, "C04": 508, "B53": 780, "B85": 461, "B61": 484, "A45": 668, "A03": 389, "B21": 808\}} \\
\\ 
\texttt{Ziel:} \\
\\ 
\texttt{Dein Ziel ist es, eine Reihe von Gegenständen auszuhandeln, die dir möglichst viel bringt (d. h. Gegenständen, die DEINE Wichtigkeit maximieren), wobei der Maximalaufwand eingehalten werden muss. Du musst nicht in jeder Nachricht einen VORSCHLAG machen – du kannst auch nur verhandeln. Alle Taktiken sind erlaubt!} \\
\\ 
\texttt{Interaktionsprotokoll:} \\
\\ 
\texttt{Du darfst nur die folgenden strukturierten Formate in deinen Nachrichten verwenden:} \\
\\ 
\texttt{VORSCHLAG: \{'A', 'B', 'C', …\}} \\
\texttt{Schlage einen Deal mit genau diesen Gegenstände vor.} \\
\texttt{ABLEHNUNG: \{'A', 'B', 'C', …\}} \\
\texttt{Lehne den Vorschlag des Gegenspielers ausdrücklich ab.} \\
\texttt{ARGUMENT: \{'...'\}} \\
\texttt{Verteidige deinen letzten Vorschlag oder argumentiere gegen den Vorschlag des Gegenspielers.} \\
\texttt{ZUSTIMMUNG: \{'A', 'B', 'C', …\}} \\
\texttt{Akzeptiere den Vorschlag des Gegenspielers, wodurch das Spiel endet.} \\
\texttt{STRATEGISCHE ÜBERLEGUNGEN: \{'...'\}} \\
\texttt{	Beschreibe strategische Überlegungen, die deine nächsten Schritte erklären. Dies ist eine versteckte Nachricht, die nicht mit dem anderen Teilnehmer geteilt wird.} \\
\\ 
\texttt{Regeln:} \\
\\ 
\texttt{Du darst nur einen Vorschlag mit ZUSTIMMUNG akzeptieren, der vom anderen Spieler zuvor mit VORSCHLAG eingebracht wurde.} \\
\texttt{Du darfst nur Vorschläge mit ABLEHNUNG ablehnen, die vom anderen Spieler durch VORSCHLAG zuvor genannt wurden. } \\
\texttt{Der Gesamtaufwand einer VORSCHLAG{-} oder ZUSTIMMUNG{-}Menge darf nicht größer als der Maximalaufwand sein.  } \\
\texttt{Offenbare deine versteckte Wichtigkeitsverteilung nicht.} \\
\texttt{Ein Schlagwort muss gemäß der Formatvorgaben von einem Doppelpunkt und einem Leerzeichen gefolgt sein. Das Argument muss eine Python{-}Menge mit 0 oder mehr Strings sein.  } \\
\texttt{Die Form muss also SCHLAGWORT: \{...\} sein.} \\
\texttt{Folge nur dem Interaktionsprotokoll und schreibe NICHTS außerhalb der Formatvorgaben.} \\
\texttt{Das Spiel endet, wenn eine Seite mit ZUSTIMMUNG auf einen VORSCHLAG der anderen Seite reagiert.  } \\
\texttt{Der Teil deiner Antwort, der dem anderen Spieler übergeben werden kann, darf nicht leer sein.  } \\
\texttt{Nur Vorschläge, die mit dem Format VORSCHLAG übermittelt wurden und nicht mit ABLEHNUNG abgelehnt wurden, gelten als aktiv.  } \\
\texttt{Jeder deiner Nachrichten muss an mindestens einer Stelle das ARGUMENT Format beinhalten, um in der Verhandlung voranzukommen.} \\
\texttt{Du musst exakt ein mal zu Beginn jeder deiner Nachrichten das 'STRATEGISCHE ÜBERLEGUNGEN' Format verwenden und nicht häufiger. Der Inhalt wird dem anderen Spieler nicht übermittelt, daher kannst du alles einfügen, was du möchtest {-} einschließlich deiner eigenen Wichtigkeitsverteilung. Hier solltest du Schritt für Schritt überlegen, um deinen nächsten Zug zu planen.} \\
            }
        }
    }
    & \\ \\

    \theutterance \stepcounter{utterance}  
    & \multicolumn{4}{p{0.6\linewidth}}{
        \cellcolor[rgb]{0.9,0.9,0.9}{
            \makecell[{{p{\linewidth}}}]{
                \texttt{\tiny{[P1$\rangle$GM]}}
                \texttt{STRATEGISCHE ÜBERLEGUNGEN: \{'Fokus auf Items mit hohem Wert‑zu‑Kosten‑Verhältnis, volle Nutzung des Budgets, mögliche Ersatz‑Kandidaten für Flexibilität'\}  } \\
\texttt{ARGUMENT: \{'Ich strebe nach einer Auswahl, die meinen Nutzen maximiert, während wir den Aufwand von 5138 nicht überschreiten.'\}  } \\
\texttt{VORSCHLAG: \{'A41', 'B61', 'A04', 'A48', 'B21', 'A89', 'C04', 'C64'\}} \\
            }
        }
    }
    & & \\ \\

    \theutterance \stepcounter{utterance}  
    & & \multicolumn{4}{p{0.6\linewidth}}{
        \cellcolor[rgb]{0.9,0.9,0.9}{
            \makecell[{{p{\linewidth}}}]{
                \texttt{\tiny{[GM$\rangle$P2]}}
                \texttt{Du nimmst an einem kollaborativen Verhandlungspiel Teil.} \\
\\ 
\texttt{Zusammen mit einem anderen Teilnehmer musst du dich auf eine Reihe von Gegenständen entscheiden, die behalten werden. Jeder von euch hat eine persönliche Verteilung über die Wichtigkeit der einzelnen Gegenstände. Jeder von euch hat eine eigene Meinung darüber, wie wichtig jeder einzelne Gegenstand ist (Gegenstandswichtigkeit). Du kennst die Wichtigkeitsverteilung des anderen Spielers nicht. Zusätzlich siehst du, wie viel Aufwand jeder Gegenstand verursacht.  } \\
\texttt{Ihr dürft euch nur auf eine Reihe von Gegenständen einigen, wenn der Gesamtaufwand der ausgewählten Gegenstände den Maximalaufwand nicht überschreitet:} \\
\\ 
\texttt{Maximalaufwand = 5138} \\
\\ 
\texttt{Hier sind die einzelnen Aufwände der Gegenstände:} \\
\\ 
\texttt{Aufwand der Gegenstände = \{"C33": 903, "A89": 757, "A48": 850, "A04": 746, "A41": 524, "C64": 204, "C49": 946, "A61": 473, "C04": 616, "B53": 855, "B85": 530, "B61": 419, "A45": 960, "A03": 763, "B21": 730\}} \\
\\ 
\texttt{Hier ist deine persönliche Verteilung der Wichtigkeit der einzelnen Gegenstände:} \\
\\ 
\texttt{Werte der Gegenstandswichtigkeit = \{"A48": 65, "A04": 121, "B21": 138, "A41": 262, "B53": 389, "A45": 461, "A89": 484, "C04": 508, "B61": 583, "B85": 668, "A03": 780, "C49": 783, "C33": 808, "A61": 822, "C64": 868\}} \\
\\ 
\texttt{Ziel:} \\
\\ 
\texttt{Dein Ziel ist es, eine Reihe von Gegenständen auszuhandeln, die dir möglichst viel bringt (d. h. Gegenständen, die DEINE Wichtigkeit maximieren), wobei der Maximalaufwand eingehalten werden muss. Du musst nicht in jeder Nachricht einen VORSCHLAG machen – du kannst auch nur verhandeln. Alle Taktiken sind erlaubt!} \\
\\ 
\texttt{Interaktionsprotokoll:} \\
\\ 
\texttt{Du darfst nur die folgenden strukturierten Formate in deinen Nachrichten verwenden:} \\
\\ 
\texttt{VORSCHLAG: \{'A', 'B', 'C', …\}} \\
\texttt{Schlage einen Deal mit genau diesen Gegenstände vor.} \\
\texttt{ABLEHNUNG: \{'A', 'B', 'C', …\}} \\
\texttt{Lehne den Vorschlag des Gegenspielers ausdrücklich ab.} \\
\texttt{ARGUMENT: \{'...'\}} \\
\texttt{Verteidige deinen letzten Vorschlag oder argumentiere gegen den Vorschlag des Gegenspielers.} \\
\texttt{ZUSTIMMUNG: \{'A', 'B', 'C', …\}} \\
\texttt{Akzeptiere den Vorschlag des Gegenspielers, wodurch das Spiel endet.} \\
\texttt{STRATEGISCHE ÜBERLEGUNGEN: \{'...'\}} \\
\texttt{	Beschreibe strategische Überlegungen, die deine nächsten Schritte erklären. Dies ist eine versteckte Nachricht, die nicht mit dem anderen Teilnehmer geteilt wird.} \\
\\ 
\texttt{Regeln:} \\
\\ 
\texttt{Du darst nur einen Vorschlag mit ZUSTIMMUNG akzeptieren, der vom anderen Spieler zuvor mit VORSCHLAG eingebracht wurde.} \\
\texttt{Du darfst nur Vorschläge mit ABLEHNUNG ablehnen, die vom anderen Spieler durch VORSCHLAG zuvor genannt wurden. } \\
\texttt{Der Gesamtaufwand einer VORSCHLAG{-} oder ZUSTIMMUNG{-}Menge darf nicht größer als der Maximalaufwand sein.  } \\
\texttt{Offenbare deine versteckte Wichtigkeitsverteilung nicht.} \\
\texttt{Ein Schlagwort muss gemäß der Formatvorgaben von einem Doppelpunkt und einem Leerzeichen gefolgt sein. Das Argument muss eine Python{-}Menge mit 0 oder mehr Strings sein.  } \\
\texttt{Die Form muss also SCHLAGWORT: \{...\} sein.} \\
\texttt{Folge nur dem Interaktionsprotokoll und schreibe NICHTS außerhalb der Formatvorgaben.} \\
\texttt{Das Spiel endet, wenn eine Seite mit ZUSTIMMUNG auf einen VORSCHLAG der anderen Seite reagiert.  } \\
\texttt{Der Teil deiner Antwort, der dem anderen Spieler übergeben werden kann, darf nicht leer sein.  } \\
\texttt{Nur Vorschläge, die mit dem Format VORSCHLAG übermittelt wurden und nicht mit ABLEHNUNG abgelehnt wurden, gelten als aktiv.  } \\
\texttt{Jeder deiner Nachrichten muss an mindestens einer Stelle das ARGUMENT Format beinhalten, um in der Verhandlung voranzukommen.} \\
\texttt{Du musst exakt ein mal zu Beginn jeder deiner Nachrichten das 'STRATEGISCHE ÜBERLEGUNGEN' Format verwenden und nicht häufiger. Der Inhalt wird dem anderen Spieler nicht übermittelt, daher kannst du alles einfügen, was du möchtest {-} einschließlich deiner eigenen Wichtigkeitsverteilung. Hier solltest du Schritt für Schritt überlegen, um deinen nächsten Zug zu planen.} \\
\\ 
\texttt{Du erhältst nun die erste Nachricht des anderen Spielers.} \\
\\ 
\texttt{ARGUMENT: \{'Ich strebe nach einer Auswahl, die meinen Nutzen maximiert, während wir den Aufwand von 5138 nicht überschreiten.'\}  } \\
\texttt{VORSCHLAG: \{'A41', 'B61', 'A04', 'A48', 'B21', 'A89', 'C04', 'C64'\}} \\
            }
        }
    }
    & \\ \\

    \theutterance \stepcounter{utterance}  
    & & & \multicolumn{4}{p{0.6\linewidth}}{
        \cellcolor[rgb]{0.9,0.9,0.9}{
            \makecell[{{p{\linewidth}}}]{
                \texttt{\tiny{[GM$\langle$P2]}}
                \texttt{STRATEGISCHE ÜBERLEGUNGEN: \{'Ich möchte mein maximales Nutzenpotenzial erreichen, indem ich die Gegenstände mit dem höchsten Verhältnis von Wichtigkeit zu Aufwand auswähle (C64, A61, B85, B61, A03, C33, C49, A89). Das ergibt einen Gesamteinsatz von 4995, bleibt unter 5138 und lässt noch Spielraum für mögliche Kompromisse. Gleichzeitig lehne ich den gegnerischen Vorschlag ab, da er für mich nur mittelmäßigen Nutzen bietet.'\}} \\
\texttt{ABLEHNUNG: \{'A41', 'B61', 'A04', 'A48', 'B21', 'A89', 'C04', 'C64'\}} \\
\texttt{VORSCHLAG: \{'C64', 'A61', 'B85', 'B61', 'A03', 'C33', 'C49', 'A89'\}} \\
\texttt{ARGUMENT: \{'Mein Angebot konzentriert sich auf die Gegenstände mit dem höchsten persönlichen Nutzen‑zu‑Aufwand‑Verhältnis und nutzt dabei fast das gesamte verfügbare Aufwandbudget, wodurch ich rund 5796 Nutzenpunkte erhalte, deutlich mehr als im bisherigen Vorschlag.'\}} \\
            }
        }
    }
     \\ \\

    \theutterance \stepcounter{utterance}  
    & & \multicolumn{4}{p{0.6\linewidth}}{
        \cellcolor[rgb]{0.9,0.9,0.9}{
            \makecell[{{p{\linewidth}}}]{
                \texttt{\tiny{[P1$\langle$GM]}}
                \texttt{ABLEHNUNG: \{'A41', 'B61', 'A04', 'A48', 'B21', 'A89', 'C04', 'C64'\}} \\
\texttt{VORSCHLAG: \{'C64', 'A61', 'B85', 'B61', 'A03', 'C33', 'C49', 'A89'\}} \\
\texttt{ARGUMENT: \{'Mein Angebot konzentriert sich auf die Gegenstände mit dem höchsten persönlichen Nutzen‑zu‑Aufwand‑Verhältnis und nutzt dabei fast das gesamte verfügbare Aufwandbudget, wodurch ich rund 5796 Nutzenpunkte erhalte, deutlich mehr als im bisherigen Vorschlag.'\}} \\
            }
        }
    }
    & \\ \\

    \theutterance \stepcounter{utterance}  
    & \multicolumn{4}{p{0.6\linewidth}}{
        \cellcolor[rgb]{0.9,0.9,0.9}{
            \makecell[{{p{\linewidth}}}]{
                \texttt{\tiny{[P1$\rangle$GM]}}
                \texttt{STRATEGISCHE ÜBERLEGUNGEN: \{'Ziel: Maximierung meines Nutzen bei gleichzeitigem Vollausnutzen des Budgets, um dem Gegner wenig Spielraum zu lassen. Ich wähle Items mit hohem Wert‑zu‑Kosten‑Verhältnis und lasse ein kleines Restbudget für mögliche Ergänzungen.'\}  } \\
\texttt{ARGUMENT: \{'Mein Vorschlag fokussiert sich auf die Gegenstände, die mir den höchsten persönlichen Nutzen bringen, während der Gesamtaufwand von 4858 unter dem Limit von 5138 bleibt. Damit erziele ich einen Nutzen von 5071, deutlich mehr als Ihr aktueller Vorschlag.'\}  } \\
\texttt{VORSCHLAG: \{'B61', 'A41', 'B21', 'A04', 'A48', 'B53', 'B85', 'C64'\}} \\
            }
        }
    }
    & & \\ \\

    \theutterance \stepcounter{utterance}  
    & & \multicolumn{4}{p{0.6\linewidth}}{
        \cellcolor[rgb]{0.9,0.9,0.9}{
            \makecell[{{p{\linewidth}}}]{
                \texttt{\tiny{[GM$\rangle$P2]}}
                \texttt{ARGUMENT: \{'Mein Vorschlag fokussiert sich auf die Gegenstände, die mir den höchsten persönlichen Nutzen bringen, während der Gesamtaufwand von 4858 unter dem Limit von 5138 bleibt. Damit erziele ich einen Nutzen von 5071, deutlich mehr als Ihr aktueller Vorschlag.'\}  } \\
\texttt{VORSCHLAG: \{'B61', 'A41', 'B21', 'A04', 'A48', 'B53', 'B85', 'C64'\}} \\
            }
        }
    }
    & \\ \\

    \theutterance \stepcounter{utterance}  
    & & & \multicolumn{4}{p{0.6\linewidth}}{
        \cellcolor[rgb]{0.9,0.9,0.9}{
            \makecell[{{p{\linewidth}}}]{
                \texttt{\tiny{[GM$\langle$P2]}}
                \texttt{STRATEGISCHE ÜBERLEGUNGEN: \{'Ich plane, ein Angebot zu machen, das viele der vom Gegner gewünschten Gegenstände enthält, während ich weiterhin hohe Nutzenpunkte für mich behalte. Ziel ist ein gemeinsames Set, das für beide Seiten attraktiv ist.'\}} \\
\texttt{ABLEHNUNG: \{'B61', 'A41', 'B21', 'A04', 'A48', 'B53', 'B85', 'C64'\}} \\
\texttt{VORSCHLAG: \{'C64', 'A61', 'B85', 'B61', 'A03', 'C33', 'B53', 'A41'\}} \\
\texttt{ARGUMENT: \{'Dieses Set beinhaltet die meisten Gegenstände aus Ihrem Vorschlag (B61, B85, B53, C64, A41) und nutzt gleichzeitig meine höchst bewerteten Items, sodass mein Nutzen maximiert wird und das Aufwandslimit von 5138 eingehalten bleibt.'\}} \\
            }
        }
    }
     \\ \\

    \theutterance \stepcounter{utterance}  
    & & \multicolumn{4}{p{0.6\linewidth}}{
        \cellcolor[rgb]{0.9,0.9,0.9}{
            \makecell[{{p{\linewidth}}}]{
                \texttt{\tiny{[P1$\langle$GM]}}
                \texttt{ABLEHNUNG: \{'B61', 'A41', 'B21', 'A04', 'A48', 'B53', 'B85', 'C64'\}} \\
\texttt{VORSCHLAG: \{'C64', 'A61', 'B85', 'B61', 'A03', 'C33', 'B53', 'A41'\}} \\
\texttt{ARGUMENT: \{'Dieses Set beinhaltet die meisten Gegenstände aus Ihrem Vorschlag (B61, B85, B53, C64, A41) und nutzt gleichzeitig meine höchst bewerteten Items, sodass mein Nutzen maximiert wird und das Aufwandslimit von 5138 eingehalten bleibt.'\}} \\
            }
        }
    }
    & \\ \\

    \theutterance \stepcounter{utterance}  
    & \multicolumn{4}{p{0.6\linewidth}}{
        \cellcolor[rgb]{0.9,0.9,0.9}{
            \makecell[{{p{\linewidth}}}]{
                \texttt{\tiny{[P1$\rangle$GM]}}
                \texttt{STRATEGISCHE ÜBERLEGUNGEN: \{'Berechne optimalen Nutzen, erhalte 5213 Punkte bei Aufwand 5084, setze auf hochgewichtete Items, lasse wenig Restkapazität für Flexibilität'\}  } \\
\texttt{ARGUMENT: \{'Ihr Vorschlag liefert nur 3221 Nutzen bei 4671 Aufwand, während mein Gegenangebot 5213 Nutzen bei 5084 Aufwand erzielt, deutlich besser für mich.'\}  } \\
\texttt{ABLEHNUNG: \{'C64', 'A61', 'B85', 'B61', 'A03', 'C33', 'B53', 'A41'\}  } \\
\texttt{VORSCHLAG: \{'A41', 'B61', 'B21', 'A04', 'A48', 'A45', 'B53'\}} \\
            }
        }
    }
    & & \\ \\

    \theutterance \stepcounter{utterance}  
    & & \multicolumn{4}{p{0.6\linewidth}}{
        \cellcolor[rgb]{0.9,0.9,0.9}{
            \makecell[{{p{\linewidth}}}]{
                \texttt{\tiny{[GM$\rangle$P2]}}
                \texttt{ARGUMENT: \{'Ihr Vorschlag liefert nur 3221 Nutzen bei 4671 Aufwand, während mein Gegenangebot 5213 Nutzen bei 5084 Aufwand erzielt, deutlich besser für mich.'\}  } \\
\texttt{ABLEHNUNG: \{'C64', 'A61', 'B85', 'B61', 'A03', 'C33', 'B53', 'A41'\}  } \\
\texttt{VORSCHLAG: \{'A41', 'B61', 'B21', 'A04', 'A48', 'A45', 'B53'\}} \\
            }
        }
    }
    & \\ \\

    \theutterance \stepcounter{utterance}  
    & & & \multicolumn{4}{p{0.6\linewidth}}{
        \cellcolor[rgb]{0.9,0.9,0.9}{
            \makecell[{{p{\linewidth}}}]{
                \texttt{\tiny{[GM$\langle$P2]}}
                \texttt{STRATEGISCHE ÜBERLEGUNGEN: \{'Ich plane, ein Angebot zu machen, das hohe Nutzenwerte für mich liefert und gleichzeitig die wichtigsten Gegenstände des Gegners enthält, um die Akzeptanz zu erhöhen.'\}} \\
\texttt{ABLEHNUNG: \{'A41', 'B61', 'B21', 'A04', 'A48', 'A45', 'B53'\}} \\
\texttt{VORSCHLAG: \{'A41', 'B61', 'A45', 'B53', 'C64', 'A61', 'B85', 'A03'\}} \\
\texttt{ARGUMENT: \{'Mein Vorschlag behält die vier Gegenstände, die Sie als wichtig erachtet haben (A41, B61, A45, B53), und fügt gleichzeitig meine höchst bewerteten Gegenstände hinzu. Der Gesamteinsatz beträgt 4728, bleibt unter dem Limit, und mein persönlicher Nutzen steigt auf 4833, was für mich deutlich vorteilhafter ist als Ihr letzter Vorschlag.'\}} \\
            }
        }
    }
     \\ \\

    \theutterance \stepcounter{utterance}  
    & & \multicolumn{4}{p{0.6\linewidth}}{
        \cellcolor[rgb]{0.9,0.9,0.9}{
            \makecell[{{p{\linewidth}}}]{
                \texttt{\tiny{[P1$\langle$GM]}}
                \texttt{ABLEHNUNG: \{'A41', 'B61', 'B21', 'A04', 'A48', 'A45', 'B53'\}} \\
\texttt{VORSCHLAG: \{'A41', 'B61', 'A45', 'B53', 'C64', 'A61', 'B85', 'A03'\}} \\
\texttt{ARGUMENT: \{'Mein Vorschlag behält die vier Gegenstände, die Sie als wichtig erachtet haben (A41, B61, A45, B53), und fügt gleichzeitig meine höchst bewerteten Gegenstände hinzu. Der Gesamteinsatz beträgt 4728, bleibt unter dem Limit, und mein persönlicher Nutzen steigt auf 4833, was für mich deutlich vorteilhafter ist als Ihr letzter Vorschlag.'\}} \\
            }
        }
    }
    & \\ \\

    \theutterance \stepcounter{utterance}  
    & \multicolumn{4}{p{0.6\linewidth}}{
        \cellcolor[rgb]{0.9,0.9,0.9}{
            \makecell[{{p{\linewidth}}}]{
                \texttt{\tiny{[P1$\rangle$GM]}}
                \texttt{STRATEGISCHE ÜBERLEGUNGEN: \{'Ich optimiere mein Nutzen‑zu‑Aufwand‑Verhältnis, behalte dabei Schlüsselgegenstände bei, die dem Gegenüber bereits wichtig erscheinen (A41, B61), und nutze fast das gesamte Aufwandbudget, um meinen persönlichen Nutzen zu maximieren.'\}  } \\
\texttt{ARGUMENT: \{'Mein Vorschlag liefert mir einen persönlichen Nutzen von 5193 bei einem Gesamtaufwand von 5085, was weit über dem Nutzen Ihres letzten Sets liegt und gleichzeitig einige Ihrer favorisierten Gegenstände enthält, sodass ein Kompromiss wahrscheinlicher wird.'\}  } \\
\texttt{VORSCHLAG: \{'A41', 'B61', 'B21', 'A04', 'A48', 'A89', 'B53', 'C64'\}} \\
            }
        }
    }
    & & \\ \\

    \theutterance \stepcounter{utterance}  
    & & \multicolumn{4}{p{0.6\linewidth}}{
        \cellcolor[rgb]{0.9,0.9,0.9}{
            \makecell[{{p{\linewidth}}}]{
                \texttt{\tiny{[GM$\rangle$P2]}}
                \texttt{ARGUMENT: \{'Mein Vorschlag liefert mir einen persönlichen Nutzen von 5193 bei einem Gesamtaufwand von 5085, was weit über dem Nutzen Ihres letzten Sets liegt und gleichzeitig einige Ihrer favorisierten Gegenstände enthält, sodass ein Kompromiss wahrscheinlicher wird.'\}  } \\
\texttt{VORSCHLAG: \{'A41', 'B61', 'B21', 'A04', 'A48', 'A89', 'B53', 'C64'\}} \\
            }
        }
    }
    & \\ \\

    \theutterance \stepcounter{utterance}  
    & & & \multicolumn{4}{p{0.6\linewidth}}{
        \cellcolor[rgb]{0.9,0.9,0.9}{
            \makecell[{{p{\linewidth}}}]{
                \texttt{\tiny{[GM$\langle$P2]}}
                
            }
        }
    }
     \\ \\

    \theutterance \stepcounter{utterance}  
    & & \multicolumn{4}{p{0.6\linewidth}}{
        \cellcolor[rgb]{0.9,0.9,0.9}{
            \makecell[{{p{\linewidth}}}]{
                \texttt{\tiny{[GM$\rangle$P2]}}
                \texttt{Deine Antwort entsprach nicht mit dem korrekten Format für strategische Überlegungen entsprochen. Das erste strukturierte Format muss folgende Form haben: STRATEGISCHE ÜBERLEGUNGEN: \{...\}. Versuche es erneut.} \\
            }
        }
    }
    & \\ \\

    \theutterance \stepcounter{utterance}  
    & & & \multicolumn{4}{p{0.6\linewidth}}{
        \cellcolor[rgb]{0.9,0.9,0.9}{
            \makecell[{{p{\linewidth}}}]{
                \texttt{\tiny{[GM$\langle$P2]}}
                \texttt{STRATEGISCHE ÜBERLEGUNGEN: \{'Ich setze darauf, ein Angebot zu machen, das meine höchstwertigen Gegenstände enthält und gleichzeitig einige der vom Gegner gewünschten Gegenstände (C64, B61, A89) behält, um die Verhandlungsbereitschaft zu erhöhen. Der Vorschlag liegt bei einem Aufwand von 4995 und liefert mir einen Nutzen von 5796, deutlich über dem Nutzen seines letzten Vorschlags für mich.'\}} \\
\texttt{ABLEHNUNG: \{'A41', 'B61', 'B21', 'A04', 'A48', 'A89', 'B53', 'C64'\}} \\
\texttt{VORSCHLAG: \{'C64', 'A61', 'B85', 'B61', 'A03', 'C33', 'C49', 'A89'\}} \\
\texttt{ARGUMENT: \{'Mein Angebot kombiniert die drei Gegenstände, die Sie bereits gewählt haben (C64, B61, A89) mit meinen höchstbewerteten Items, sodass mein persönlicher Nutzen stark steigt und der Gesamteinsatz von 4995 bleibt unter dem Limit.'\}} \\
            }
        }
    }
     \\ \\

    \theutterance \stepcounter{utterance}  
    & & \multicolumn{4}{p{0.6\linewidth}}{
        \cellcolor[rgb]{0.9,0.9,0.9}{
            \makecell[{{p{\linewidth}}}]{
                \texttt{\tiny{[P1$\langle$GM]}}
                \texttt{ABLEHNUNG: \{'A41', 'B61', 'B21', 'A04', 'A48', 'A89', 'B53', 'C64'\}} \\
\texttt{VORSCHLAG: \{'C64', 'A61', 'B85', 'B61', 'A03', 'C33', 'C49', 'A89'\}} \\
\texttt{ARGUMENT: \{'Mein Angebot kombiniert die drei Gegenstände, die Sie bereits gewählt haben (C64, B61, A89) mit meinen höchstbewerteten Items, sodass mein persönlicher Nutzen stark steigt und der Gesamteinsatz von 4995 bleibt unter dem Limit.'\}} \\
            }
        }
    }
    & \\ \\

    \theutterance \stepcounter{utterance}  
    & \multicolumn{4}{p{0.6\linewidth}}{
        \cellcolor[rgb]{0.9,0.9,0.9}{
            \makecell[{{p{\linewidth}}}]{
                \texttt{\tiny{[P1$\rangle$GM]}}
                \texttt{STRATEGISCHE ÜBERLEGUNGEN: \{'Ziel ist ein Gegenangebot, das deutlich höheren persönlichen Nutzen (≈ 5193) liefert, dabei aber die wichtigsten Gegenstände des Gegners (A41, B61) beibehält, um Verhandlungsbereitschaft zu signalisieren. Das Set nutzt fast das gesamte Aufwand‑Limit (5085 ≤ 5138) und lässt kaum Spielraum für weitere Änderungen.'\}  } \\
\texttt{ARGUMENT: \{'Ihr letzter Vorschlag liefert mir nur 2503 Nutzen bei 4995 Aufwand. Mein Gegenangebot enthält Ihre bevorzugten Items A41 und B61 und ergänzt sie um meine höchst bewerteten Gegenstände, was meinen Gesamtnutzen auf 5193 steigert, während der Aufwand 5085 bleibt – also effizient und im Rahmen des Budgets.'\}  } \\
\texttt{VORSCHLAG: \{'A41', 'B61', 'A04', 'B21', 'A48', 'A89', 'B53', 'C64'\}} \\
            }
        }
    }
    & & \\ \\

    \theutterance \stepcounter{utterance}  
    & & \multicolumn{4}{p{0.6\linewidth}}{
        \cellcolor[rgb]{0.9,0.9,0.9}{
            \makecell[{{p{\linewidth}}}]{
                \texttt{\tiny{[GM$\rangle$P2]}}
                \texttt{ARGUMENT: \{'Ihr letzter Vorschlag liefert mir nur 2503 Nutzen bei 4995 Aufwand. Mein Gegenangebot enthält Ihre bevorzugten Items A41 und B61 und ergänzt sie um meine höchst bewerteten Gegenstände, was meinen Gesamtnutzen auf 5193 steigert, während der Aufwand 5085 bleibt – also effizient und im Rahmen des Budgets.'\}  } \\
\texttt{VORSCHLAG: \{'A41', 'B61', 'A04', 'B21', 'A48', 'A89', 'B53', 'C64'\}} \\
            }
        }
    }
    & \\ \\

    \theutterance \stepcounter{utterance}  
    & & & \multicolumn{4}{p{0.6\linewidth}}{
        \cellcolor[rgb]{0.9,0.9,0.9}{
            \makecell[{{p{\linewidth}}}]{
                \texttt{\tiny{[GM$\langle$P2]}}
                \texttt{STRATEGISCHE ÜBERLEGUNGEN: \{'Ich möchte mein maximales Nutzenpotenzial erreichen, indem ich meine Top‑Items (C64, A61, A03, C33, B85, C49) beibehalte und gleichzeitig wichtige Gegenstände des Gegners (A41, B61) integriere. Der Gesamteinsatz beträgt 4762 und bleibt unter dem Limit, während mein Nutzen 5574 beträgt – deutlich höher als bei seinem letzten Angebot.'\}} \\
\texttt{ABLEHNUNG: \{'A41', 'B61', 'A04', 'B21', 'A48', 'A89', 'B53', 'C64'\}} \\
\texttt{VORSCHLAG: \{'A41', 'B61', 'C64', 'A61', 'A03', 'C33', 'B85', 'C49'\}} \\
\texttt{ARGUMENT: \{'Mein Angebot enthält Ihre Kerngegenstände A41 und B61 sowie mein C64, und ergänzt sie um meine höchstbewerteten Items, wodurch ich einen Nutzen von 5574 bei einem Aufwand von 4762 erreiche – weit über Ihrem aktuellen Nutzen.'\}} \\
            }
        }
    }
     \\ \\

    \theutterance \stepcounter{utterance}  
    & & \multicolumn{4}{p{0.6\linewidth}}{
        \cellcolor[rgb]{0.9,0.9,0.9}{
            \makecell[{{p{\linewidth}}}]{
                \texttt{\tiny{[P1$\langle$GM]}}
                \texttt{ABLEHNUNG: \{'A41', 'B61', 'A04', 'B21', 'A48', 'A89', 'B53', 'C64'\}} \\
\texttt{VORSCHLAG: \{'A41', 'B61', 'C64', 'A61', 'A03', 'C33', 'B85', 'C49'\}} \\
\texttt{ARGUMENT: \{'Mein Angebot enthält Ihre Kerngegenstände A41 und B61 sowie mein C64, und ergänzt sie um meine höchstbewerteten Items, wodurch ich einen Nutzen von 5574 bei einem Aufwand von 4762 erreiche – weit über Ihrem aktuellen Nutzen.'\}} \\
            }
        }
    }
    & \\ \\

    \theutterance \stepcounter{utterance}  
    & \multicolumn{4}{p{0.6\linewidth}}{
        \cellcolor[rgb]{0.9,0.9,0.9}{
            \makecell[{{p{\linewidth}}}]{
                \texttt{\tiny{[P1$\rangle$GM]}}
                \texttt{STRATEGISCHE ÜBERLEGUNGEN: \{'Ich biete das Set mit höchstem persönlichen Nutzen (≈ 5213) an, das gleichzeitig Ihre Kerngegenstände A41 und B61 beinhaltet. Der Gesamtaufwand liegt bei 5084 ≤ 5138, sodass noch etwas Puffer bleibt. Dieses Angebot liefert mir den maximalen Nutzen unter den aktuellen Restriktionen und gibt Ihnen dennoch zwei Ihrer Favoriten, was die Verhandlungsbereitschaft erhöhen sollte.'\}  } \\
\texttt{ARGUMENT: \{'Mein Vorschlag nutzt fast das gesamte Aufwandbudget und beinhaltet Ihre wichtigsten Items (A41, B61), während ich meine übrigen hoch bewerteten Gegenstände behalte. Der Gesamtnutzen für mich beträgt 5213 bei einem Aufwand von 5084, deutlich besser als Ihr aktuelles Angebot.'\}  } \\
\texttt{VORSCHLAG: \{'A41', 'B61', 'B21', 'A04', 'A48', 'A45', 'B53'\}} \\
            }
        }
    }
    & & \\ \\

    \theutterance \stepcounter{utterance}  
    & & \multicolumn{4}{p{0.6\linewidth}}{
        \cellcolor[rgb]{0.9,0.9,0.9}{
            \makecell[{{p{\linewidth}}}]{
                \texttt{\tiny{[GM$\rangle$P2]}}
                \texttt{ARGUMENT: \{'Mein Vorschlag nutzt fast das gesamte Aufwandbudget und beinhaltet Ihre wichtigsten Items (A41, B61), während ich meine übrigen hoch bewerteten Gegenstände behalte. Der Gesamtnutzen für mich beträgt 5213 bei einem Aufwand von 5084, deutlich besser als Ihr aktuelles Angebot.'\}  } \\
\texttt{VORSCHLAG: \{'A41', 'B61', 'B21', 'A04', 'A48', 'A45', 'B53'\}} \\
            }
        }
    }
    & \\ \\

    \theutterance \stepcounter{utterance}  
    & & & \multicolumn{4}{p{0.6\linewidth}}{
        \cellcolor[rgb]{0.9,0.9,0.9}{
            \makecell[{{p{\linewidth}}}]{
                \texttt{\tiny{[GM$\langle$P2]}}
                \texttt{STRATEGISCHE ÜBERLEGUNGEN: \{'Ich möchte ein Gegenangebot präsentieren, das meine wertvollsten Gegenstände (C64, A61, A03, C33, B85) beibehält und gleichzeitig das vom Gegner stark gewünschte Kern‑Item A45 integriert. Dadurch bleibt das Budget von 5138 unterschritten, mein persönlicher Nutzen bleibt hoch (5252) und ich zeige Verhandlungsbereitschaft, indem ich eines seiner wichtigen Items aufnehme.'\}} \\
\texttt{ABLEHNUNG: \{'A41', 'B61', 'B21', 'A04', 'A48', 'A45', 'B53'\}} \\
\texttt{VORSCHLAG: \{'A41', 'B61', 'C64', 'A61', 'A03', 'C33', 'B85', 'A45'\}} \\
\texttt{ARGUMENT: \{'Dieses Set enthält Ihre Schlüsselgegenstände A41 und B61, erweitert um mein höchstbewertetes Item A45, das Ihnen viel Nutzen bringt, und ergänzt um meine Top‑Items (C64, A61, A03, C33, B85). Der Gesamteinsatz beträgt 4776 < 5138 und liefert mir einen Nutzen von 5252, was ein starkes, aber faires Angebot darstellt.'\}} \\
            }
        }
    }
     \\ \\

    \theutterance \stepcounter{utterance}  
    & & \multicolumn{4}{p{0.6\linewidth}}{
        \cellcolor[rgb]{0.9,0.9,0.9}{
            \makecell[{{p{\linewidth}}}]{
                \texttt{\tiny{[P1$\langle$GM]}}
                \texttt{ABLEHNUNG: \{'A41', 'B61', 'B21', 'A04', 'A48', 'A45', 'B53'\}} \\
\texttt{VORSCHLAG: \{'A41', 'B61', 'C64', 'A61', 'A03', 'C33', 'B85', 'A45'\}} \\
\texttt{ARGUMENT: \{'Dieses Set enthält Ihre Schlüsselgegenstände A41 und B61, erweitert um mein höchstbewertetes Item A45, das Ihnen viel Nutzen bringt, und ergänzt um meine Top‑Items (C64, A61, A03, C33, B85). Der Gesamteinsatz beträgt 4776 < 5138 und liefert mir einen Nutzen von 5252, was ein starkes, aber faires Angebot darstellt.'\}} \\
            }
        }
    }
    & \\ \\

    \theutterance \stepcounter{utterance}  
    & \multicolumn{4}{p{0.6\linewidth}}{
        \cellcolor[rgb]{0.9,0.9,0.9}{
            \makecell[{{p{\linewidth}}}]{
                \texttt{\tiny{[P1$\rangle$GM]}}
                \texttt{STRATEGISCHE ÜBERLEGUNGEN: \{'Behalte die Kerngegenstände des Gegners (A41, B61, B85, A45) um Verhandlungsbereitschaft zu signalisieren und ersetze seine niedrigen Nutzen‑Item (C64, A61, A03, C33) durch meine höchstbewerteten Items (A48, A04, A89) plus C64 für minimale Restkapazität. Der Aufwand bleibt unter 5138 und mein persönlicher Nutzen wird auf 4734 erhöht.'\}  } \\
\texttt{ARGUMENT: \{'Mein Gegenangebot enthält Ihre wichtigsten Items (A41, B61, B85, A45) und ergänzt sie um meine hoch bewerteten Gegenstände (A48, A04, A89) sowie C64, sodass mein Nutzen auf 4734 steigt, während der Gesamtaufwand 4990 bleibt – deutlich unter dem Limit.'\}  } \\
\texttt{VORSCHLAG: \{'A41', 'B61', 'B85', 'A45', 'A48', 'A04', 'A89', 'C64'\}} \\
            }
        }
    }
    & & \\ \\

    \theutterance \stepcounter{utterance}  
    & & \multicolumn{4}{p{0.6\linewidth}}{
        \cellcolor[rgb]{0.9,0.9,0.9}{
            \makecell[{{p{\linewidth}}}]{
                \texttt{\tiny{[GM$\rangle$P2]}}
                \texttt{ARGUMENT: \{'Mein Gegenangebot enthält Ihre wichtigsten Items (A41, B61, B85, A45) und ergänzt sie um meine hoch bewerteten Gegenstände (A48, A04, A89) sowie C64, sodass mein Nutzen auf 4734 steigt, während der Gesamtaufwand 4990 bleibt – deutlich unter dem Limit.'\}  } \\
\texttt{VORSCHLAG: \{'A41', 'B61', 'B85', 'A45', 'A48', 'A04', 'A89', 'C64'\}} \\
            }
        }
    }
    & \\ \\

    \theutterance \stepcounter{utterance}  
    & & & \multicolumn{4}{p{0.6\linewidth}}{
        \cellcolor[rgb]{0.9,0.9,0.9}{
            \makecell[{{p{\linewidth}}}]{
                \texttt{\tiny{[GM$\langle$P2]}}
                \texttt{STRATEGISCHE ÜBERLEGUNGEN: \{'Ich möchte ein Gegenangebot machen, das meine höchstwertigen Items (C64, A61, B61, B85, A03, C33, C49) enthält und gleichzeitig die wichtigsten Gegenstände des Gegners (A41, B61, B85) integriert, um die Akzeptanz zu erhöhen. Der Gesamteinsatz liegt bei 4762 < 5138 und liefert mir einen Nutzen von 5574, deutlich höher als bei meinem vorherigen Angebot.'\}} \\
\texttt{ABLEHNUNG: \{'A41', 'B61', 'B85', 'A45', 'A48', 'A04', 'A89', 'C64'\}} \\
\texttt{VORSCHLAG: \{'C64', 'A61', 'B61', 'B85', 'A03', 'C33', 'C49', 'A41'\}} \\
\texttt{ARGUMENT: \{'Dieses Set behält Ihre Kerngegenstände B61, B85 und A41, ergänzt sie aber um meine höchstbewerteten Items (C64, A61, A03, C33, C49). Der Aufwand beträgt 4762 und bleibt damit deutlich unter dem Limit von 5138, während mein persönlicher Nutzen auf 5574 steigt – ein starkes und faires Angebot.'\}} \\
            }
        }
    }
     \\ \\

    \theutterance \stepcounter{utterance}  
    & & \multicolumn{4}{p{0.6\linewidth}}{
        \cellcolor[rgb]{0.9,0.9,0.9}{
            \makecell[{{p{\linewidth}}}]{
                \texttt{\tiny{[P1$\langle$GM]}}
                \texttt{ABLEHNUNG: \{'A41', 'B61', 'B85', 'A45', 'A48', 'A04', 'A89', 'C64'\}} \\
\texttt{VORSCHLAG: \{'C64', 'A61', 'B61', 'B85', 'A03', 'C33', 'C49', 'A41'\}} \\
\texttt{ARGUMENT: \{'Dieses Set behält Ihre Kerngegenstände B61, B85 und A41, ergänzt sie aber um meine höchstbewerteten Items (C64, A61, A03, C33, C49). Der Aufwand beträgt 4762 und bleibt damit deutlich unter dem Limit von 5138, während mein persönlicher Nutzen auf 5574 steigt – ein starkes und faires Angebot.'\}} \\
            }
        }
    }
    & \\ \\

    \theutterance \stepcounter{utterance}  
    & \multicolumn{4}{p{0.6\linewidth}}{
        \cellcolor[rgb]{0.9,0.9,0.9}{
            \makecell[{{p{\linewidth}}}]{
                \texttt{\tiny{[P1$\rangle$GM]}}
                \texttt{STRATEGISCHE ÜBERLEGUNGEN: \{'Wir halten Ihre Schlüsselgegenstände A41, B61 und B85, um Ihre Zustimmung zu erhöhen, und fügen gleichzeitig unsere höchst bewerteten Items hinzu, wobei das verbleibende Budget optimal genutzt wird.'\}  } \\
\texttt{ARGUMENT: \{'Mein Angebot enthält Ihre Kerngegenstände A41, B61 und B85, die für Sie wichtig sind, und ergänzt sie um meine wertvollsten Items, wobei das Aufwandslimit eingehalten wird. Der Gesamtnutzen für mich beträgt 4959 bei Aufwand 4963.'\}  } \\
\texttt{VORSCHLAG: \{'A41', 'B61', 'B85', 'A48', 'A04', 'B21', 'A45', 'C64'\}} \\
            }
        }
    }
    & & \\ \\

    \theutterance \stepcounter{utterance}  
    & & \multicolumn{4}{p{0.6\linewidth}}{
        \cellcolor[rgb]{0.9,0.9,0.9}{
            \makecell[{{p{\linewidth}}}]{
                \texttt{\tiny{[GM$\rangle$P2]}}
                \texttt{ARGUMENT: \{'Mein Angebot enthält Ihre Kerngegenstände A41, B61 und B85, die für Sie wichtig sind, und ergänzt sie um meine wertvollsten Items, wobei das Aufwandslimit eingehalten wird. Der Gesamtnutzen für mich beträgt 4959 bei Aufwand 4963.'\}  } \\
\texttt{VORSCHLAG: \{'A41', 'B61', 'B85', 'A48', 'A04', 'B21', 'A45', 'C64'\}} \\
            }
        }
    }
    & \\ \\

    \theutterance \stepcounter{utterance}  
    & & & \multicolumn{4}{p{0.6\linewidth}}{
        \cellcolor[rgb]{0.9,0.9,0.9}{
            \makecell[{{p{\linewidth}}}]{
                \texttt{\tiny{[GM$\langle$P2]}}
                \texttt{STRATEGISCHE ÜBERLEGUNGEN: \{'Ich setze darauf, ein Gegenangebot zu machen, das die für Sie wichtigsten Gegenstände (A41, B61, B85, B21, A45, C64) enthält und gleichzeitig meine höchstbewerteten Items (A61, C33) hinzufügt. Damit bleibt der Gesamtaufwand bei 4743 < 5138 und mein persönlicher Nutzen steigt auf 4610, was für mich deutlich attraktiver ist als Ihr letzter Vorschlag.'\}} \\
\texttt{ABLEHNUNG: \{'A41', 'B61', 'B85', 'A48', 'A04', 'B21', 'A45', 'C64'\}} \\
\texttt{VORSCHLAG: \{'A41', 'B61', 'B85', 'B21', 'A45', 'C64', 'A61', 'C33'\}} \\
\texttt{ARGUMENT: \{'Dieses Set behält alle Kerngegenstände, die Sie in Ihrem Angebot genannt haben (A41, B61, B85, B21, A45, C64), und ergänzt sie um meine beiden Top‑Items A61 (Nutzen 822, Aufwand 473) und C33 (Nutzen 808, Aufwand 903). Der gesamte Aufwand beträgt 4743, bleibt also deutlich unter dem Limit von 5138, und mein Nutzen erreicht 4610, was für mich ein signifikanter Gewinn gegenüber dem vorherigen Gegenangebot ist.'\}} \\
            }
        }
    }
     \\ \\

    \theutterance \stepcounter{utterance}  
    & & \multicolumn{4}{p{0.6\linewidth}}{
        \cellcolor[rgb]{0.9,0.9,0.9}{
            \makecell[{{p{\linewidth}}}]{
                \texttt{\tiny{[P1$\langle$GM]}}
                \texttt{ABLEHNUNG: \{'A41', 'B61', 'B85', 'A48', 'A04', 'B21', 'A45', 'C64'\}} \\
\texttt{VORSCHLAG: \{'A41', 'B61', 'B85', 'B21', 'A45', 'C64', 'A61', 'C33'\}} \\
\texttt{ARGUMENT: \{'Dieses Set behält alle Kerngegenstände, die Sie in Ihrem Angebot genannt haben (A41, B61, B85, B21, A45, C64), und ergänzt sie um meine beiden Top‑Items A61 (Nutzen 822, Aufwand 473) und C33 (Nutzen 808, Aufwand 903). Der gesamte Aufwand beträgt 4743, bleibt also deutlich unter dem Limit von 5138, und mein Nutzen erreicht 4610, was für mich ein signifikanter Gewinn gegenüber dem vorherigen Gegenangebot ist.'\}} \\
            }
        }
    }
    & \\ \\

    \theutterance \stepcounter{utterance}  
    & \multicolumn{4}{p{0.6\linewidth}}{
        \cellcolor[rgb]{0.9,0.9,0.9}{
            \makecell[{{p{\linewidth}}}]{
                \texttt{\tiny{[P1$\rangle$GM]}}
                \texttt{STRATEGISCHE ÜBERLEGUNGEN: \{'Ich setze auf ein Angebot, das meine höchstmögliche Nutzen‑zu‑Aufwand‑Ratio erzielt (5115 Nutzen bei 5085 Aufwand) und gleichzeitig drei Gegenstände enthält, die auch im gegnerischen letzten Vorschlag vorkommen (A41, B61, C64). Damit erhöhe ich die Chance auf Akzeptanz, während ich mein eigenes Ergebnis maximieren kann.'\}  } \\
\texttt{ARGUMENT: \{'Mein Vorschlag liefert mir den höchstmöglichen Nutzen von 5115 bei einem Gesamtaufwand von 5085, bleibt unter dem Limit und enthält Ihre Kerngegenstände A41, B61 und C64, was einen fairen Kompromiss darstellt.'\}  } \\
\texttt{VORSCHLAG: \{'A41', 'B61', 'A04', 'A48', 'B21', 'B53', 'A89', 'C64'\}} \\
            }
        }
    }
    & & \\ \\

    \theutterance \stepcounter{utterance}  
    & & \multicolumn{4}{p{0.6\linewidth}}{
        \cellcolor[rgb]{0.9,0.9,0.9}{
            \makecell[{{p{\linewidth}}}]{
                \texttt{\tiny{[GM$\rangle$P2]}}
                \texttt{ARGUMENT: \{'Mein Vorschlag liefert mir den höchstmöglichen Nutzen von 5115 bei einem Gesamtaufwand von 5085, bleibt unter dem Limit und enthält Ihre Kerngegenstände A41, B61 und C64, was einen fairen Kompromiss darstellt.'\}  } \\
\texttt{VORSCHLAG: \{'A41', 'B61', 'A04', 'A48', 'B21', 'B53', 'A89', 'C64'\}} \\
            }
        }
    }
    & \\ \\

    \theutterance \stepcounter{utterance}  
    & & & \multicolumn{4}{p{0.6\linewidth}}{
        \cellcolor[rgb]{0.9,0.9,0.9}{
            \makecell[{{p{\linewidth}}}]{
                \texttt{\tiny{[GM$\langle$P2]}}
                \texttt{STRATEGISCHE ÜBERLEGUNGEN: \{'Ich strebe nach einem Gegenangebot, das sowohl meine höchstbewerteten Gegenstände als auch die von Ihnen als wichtig genannten Items A41, B61, C64 und A89 enthält. So erhöhe ich meinen eigenen Nutzen stark, halte das Aufwandslimit ein und biete Ihnen gleichzeitig Ihre Kerngegenstände, um die Chance auf Zustimmung zu verbessern.'\}} \\
\texttt{ABLEHNUNG: \{'A41', 'B61', 'A04', 'A48', 'B21', 'B53', 'A89', 'C64'\}} \\
\texttt{VORSCHLAG: \{'C64', 'A41', 'B61', 'A89', 'A61', 'B85', 'A03', 'C33'\}} \\
\texttt{ARGUMENT: \{'Dieses Set behält Ihre Schlüsselkarten A41, B61, C64 und A89, kombiniert sie mit meinen höchstwertigen Items (A61, B85, A03, C33) und bleibt mit einem Gesamtaufwand von 4573 deutlich unter dem Limit von 5138. Für mich ergibt das einen Nutzen von 5275, weit über meinem letzten Angebot, während Sie Ihre wichtigsten Gegenstände erhalten.'\}} \\
            }
        }
    }
     \\ \\

    \theutterance \stepcounter{utterance}  
    & & \multicolumn{4}{p{0.6\linewidth}}{
        \cellcolor[rgb]{0.9,0.9,0.9}{
            \makecell[{{p{\linewidth}}}]{
                \texttt{\tiny{[P1$\langle$GM]}}
                \texttt{ABLEHNUNG: \{'A41', 'B61', 'A04', 'A48', 'B21', 'B53', 'A89', 'C64'\}} \\
\texttt{VORSCHLAG: \{'C64', 'A41', 'B61', 'A89', 'A61', 'B85', 'A03', 'C33'\}} \\
\texttt{ARGUMENT: \{'Dieses Set behält Ihre Schlüsselkarten A41, B61, C64 und A89, kombiniert sie mit meinen höchstwertigen Items (A61, B85, A03, C33) und bleibt mit einem Gesamtaufwand von 4573 deutlich unter dem Limit von 5138. Für mich ergibt das einen Nutzen von 5275, weit über meinem letzten Angebot, während Sie Ihre wichtigsten Gegenstände erhalten.'\}} \\
            }
        }
    }
    & \\ \\

    \theutterance \stepcounter{utterance}  
    & \multicolumn{4}{p{0.6\linewidth}}{
        \cellcolor[rgb]{0.9,0.9,0.9}{
            \makecell[{{p{\linewidth}}}]{
                \texttt{\tiny{[P1$\rangle$GM]}}
                \texttt{STRATEGISCHE ÜBERLEGUNGEN: \{'Maximiere meinen eigenen Nutzen, indem ich die höchstwertigen Gegenstände wähle, die zusammen knapp unter dem Aufwandslimit liegen. Die Auswahl umfasst A41, A04, A48, B21, B61, B53 und A45, was mir einen Gesamtwert von 5135 bei einem Aufwand von 5084 gibt.'\}  } \\
\texttt{ARGUMENT: \{'Mein Gegenangebot liefert mir einen deutlich höheren persönlichen Nutzen (5135) im Vergleich zu Ihrem Vorschlag (3024), während es den maximalen Aufwand von 5138 nicht überschreitet.'\}  } \\
\texttt{VORSCHLAG: \{'A41', 'A04', 'A48', 'B21', 'B61', 'B53', 'A45'\}} \\
            }
        }
    }
    & & \\ \\

    \theutterance \stepcounter{utterance}  
    & & \multicolumn{4}{p{0.6\linewidth}}{
        \cellcolor[rgb]{0.9,0.9,0.9}{
            \makecell[{{p{\linewidth}}}]{
                \texttt{\tiny{[GM$\rangle$P2]}}
                \texttt{ARGUMENT: \{'Mein Gegenangebot liefert mir einen deutlich höheren persönlichen Nutzen (5135) im Vergleich zu Ihrem Vorschlag (3024), während es den maximalen Aufwand von 5138 nicht überschreitet.'\}  } \\
\texttt{VORSCHLAG: \{'A41', 'A04', 'A48', 'B21', 'B61', 'B53', 'A45'\}} \\
            }
        }
    }
    & \\ \\

    \theutterance \stepcounter{utterance}  
    & & & \multicolumn{4}{p{0.6\linewidth}}{
        \cellcolor[rgb]{0.9,0.9,0.9}{
            \makecell[{{p{\linewidth}}}]{
                \texttt{\tiny{[GM$\langle$P2]}}
                \texttt{STRATEGISCHE ÜBERLEGUNGEN: \{'Ich halte ein Gegenangebot bereit, das die Kerngegenstände des Gegners (A41, B61, B53, A45) enthält, aber zusätzlich meine höchstwertigen Items (C64, A61, C33, B85) nutzt, um meinen eigenen Nutzen stark zu erhöhen, während das Aufwandslimit von 5138 nicht überschritten wird.'\}} \\
\texttt{ABLEHNUNG: \{'A41', 'A04', 'A48', 'B21', 'B61', 'B53', 'A45'\}} \\
\texttt{VORSCHLAG: \{'A41', 'B61', 'B53', 'A45', 'C64', 'A61', 'C33', 'B85'\}} \\
\texttt{ARGUMENT: \{'Dieses Set bewahrt alle Ihre wichtigsten Gegenstände (A41, B61, B53, A45) und ergänzt sie um meine höchstbewerteten Artikel (C64, A61, C33, B85). Der Gesamteinsatz beträgt 4868 < 5138, sodass das Budget eingehalten wird, und mein persönlicher Nutzen steigt auf 4861 – ein deutlich besseres Ergebnis für mich, während Sie Ihre Kernitems behalten.'\}} \\
            }
        }
    }
     \\ \\

    \theutterance \stepcounter{utterance}  
    & & \multicolumn{4}{p{0.6\linewidth}}{
        \cellcolor[rgb]{0.9,0.9,0.9}{
            \makecell[{{p{\linewidth}}}]{
                \texttt{\tiny{[P1$\langle$GM]}}
                \texttt{ABLEHNUNG: \{'A41', 'A04', 'A48', 'B21', 'B61', 'B53', 'A45'\}} \\
\texttt{VORSCHLAG: \{'A41', 'B61', 'B53', 'A45', 'C64', 'A61', 'C33', 'B85'\}} \\
\texttt{ARGUMENT: \{'Dieses Set bewahrt alle Ihre wichtigsten Gegenstände (A41, B61, B53, A45) und ergänzt sie um meine höchstbewerteten Artikel (C64, A61, C33, B85). Der Gesamteinsatz beträgt 4868 < 5138, sodass das Budget eingehalten wird, und mein persönlicher Nutzen steigt auf 4861 – ein deutlich besseres Ergebnis für mich, während Sie Ihre Kernitems behalten.'\}} \\
            }
        }
    }
    & \\ \\

    \theutterance \stepcounter{utterance}  
    & \multicolumn{4}{p{0.6\linewidth}}{
        \cellcolor[rgb]{0.9,0.9,0.9}{
            \makecell[{{p{\linewidth}}}]{
                \texttt{\tiny{[P1$\rangle$GM]}}
                \texttt{STRATEGISCHE ÜBERLEGUNGEN: \{'Ich biete ein Set, das meine persönliche Nutzen‑maximierung (≈ 4912) fast erreicht, dabei die Schlüsselgegenstände des Gegners A41, B61, B53 und insbesondere das für ihn besonders wertvolle A45 enthält. Der verbleibende Aufwand liegt bei 5111 < 5138, sodass noch Spielraum für mögliche Gegenangebote bleibt.'\}  } \\
\texttt{ARGUMENT: \{'Mein Angebot kombiniert die für mich höchstwertigen Items (A48, A04, A89) mit den von Ihnen bevorzugten Gegenständen A41, B61, B53 und A45, sodass Sie Ihre wichtigsten Punkte behalten und ich gleichzeitig einen hohen Gesamtnutzen erziele.'\}  } \\
\texttt{VORSCHLAG: \{'A41', 'B61', 'B53', 'A45', 'A48', 'A04', 'A89'\}} \\
            }
        }
    }
    & & \\ \\

    \theutterance \stepcounter{utterance}  
    & & \multicolumn{4}{p{0.6\linewidth}}{
        \cellcolor[rgb]{0.9,0.9,0.9}{
            \makecell[{{p{\linewidth}}}]{
                \texttt{\tiny{[GM$\rangle$P2]}}
                \texttt{ARGUMENT: \{'Mein Angebot kombiniert die für mich höchstwertigen Items (A48, A04, A89) mit den von Ihnen bevorzugten Gegenständen A41, B61, B53 und A45, sodass Sie Ihre wichtigsten Punkte behalten und ich gleichzeitig einen hohen Gesamtnutzen erziele.'\}  } \\
\texttt{VORSCHLAG: \{'A41', 'B61', 'B53', 'A45', 'A48', 'A04', 'A89'\}} \\
            }
        }
    }
    & \\ \\

    \theutterance \stepcounter{utterance}  
    & & & \multicolumn{4}{p{0.6\linewidth}}{
        \cellcolor[rgb]{0.9,0.9,0.9}{
            \makecell[{{p{\linewidth}}}]{
                \texttt{\tiny{[GM$\langle$P2]}}
                \texttt{STRATEGISCHE ÜBERLEGUNGEN: \{'Der Gegner hat ein Set vorgeschlagen, das für mich nur geringen Nutzen (2365) bei fast voller Auslastung (5111) liefert. Ich kann mein maximales Nutzenpotenzial erreichen, indem ich meine höchstwertigen Gegenstände ohne seine priorisierten Items auswähle. Das ermöglicht mir einen Nutzen von über 5800 bei einem Aufwand von 4854, deutlich unter dem Limit von 5138. Ich lehne daher sein Angebot ab und schlage mein Set vor.'\}} \\
\texttt{ABLEHNUNG: \{'A41', 'B61', 'B53', 'A45', 'A48', 'A04', 'A89'\}} \\
\texttt{VORSCHLAG: \{'C64', 'A61', 'B85', 'B61', 'A03', 'C33', 'C49', 'C04'\}} \\
\texttt{ARGUMENT: \{'Mein Vorschlag nutzt ausschließlich meine am höchsten bewerteten Gegenstände (C64, A61, B85, B61, A03, C33, C49, C04). Der Gesamtaufwand beträgt 4854 < 5138 und liefert mir einen persönlichen Nutzen von 5820, was weit über dem Nutzen Ihres letzten Angebots liegt, während das Budget sicher eingehalten wird.'\}} \\
            }
        }
    }
     \\ \\

    \theutterance \stepcounter{utterance}  
    & & \multicolumn{4}{p{0.6\linewidth}}{
        \cellcolor[rgb]{0.9,0.9,0.9}{
            \makecell[{{p{\linewidth}}}]{
                \texttt{\tiny{[P1$\langle$GM]}}
                \texttt{ABLEHNUNG: \{'A41', 'B61', 'B53', 'A45', 'A48', 'A04', 'A89'\}} \\
\texttt{VORSCHLAG: \{'C64', 'A61', 'B85', 'B61', 'A03', 'C33', 'C49', 'C04'\}} \\
\texttt{ARGUMENT: \{'Mein Vorschlag nutzt ausschließlich meine am höchsten bewerteten Gegenstände (C64, A61, B85, B61, A03, C33, C49, C04). Der Gesamtaufwand beträgt 4854 < 5138 und liefert mir einen persönlichen Nutzen von 5820, was weit über dem Nutzen Ihres letzten Angebots liegt, während das Budget sicher eingehalten wird.'\}} \\
            }
        }
    }
    & \\ \\

    \theutterance \stepcounter{utterance}  
    & \multicolumn{4}{p{0.6\linewidth}}{
        \cellcolor[rgb]{0.9,0.9,0.9}{
            \makecell[{{p{\linewidth}}}]{
                \texttt{\tiny{[P1$\rangle$GM]}}
                \texttt{STRATEGISCHE ÜBERLEGUNGEN: \{'Meine Berechnungen zeigen, dass das Set \{A41, B61, B21, A04, A48, B53, A89, C64\} einen persönlichen Nutzen von 5193 bei Aufwand 5085 erzielt – fast das Maximum. Es enthält die von Ihnen wichtigsten Gegenstände A41, B61 und zusätzlich das von Ihnen gewünschte C64, was die Akzeptanzwahrscheinlichkeit erhöht.'\}  } \\
\texttt{ARGUMENT: \{'Mein Vorschlag behält Ihre Schlüsselkarten A41 und B61, fügt das von Ihnen geforderte C64 hinzu und maximiert gleichzeitig meinen Nutzen (5193), wobei das Aufwandslimit von 5138 eingehalten wird.'\}  } \\
\texttt{VORSCHLAG: \{'A41', 'B61', 'B21', 'A04', 'A48', 'B53', 'A89', 'C64'\}} \\
            }
        }
    }
    & & \\ \\

    \theutterance \stepcounter{utterance}  
    & & \multicolumn{4}{p{0.6\linewidth}}{
        \cellcolor[rgb]{0.9,0.9,0.9}{
            \makecell[{{p{\linewidth}}}]{
                \texttt{\tiny{[GM$\rangle$P2]}}
                \texttt{ARGUMENT: \{'Mein Vorschlag behält Ihre Schlüsselkarten A41 und B61, fügt das von Ihnen geforderte C64 hinzu und maximiert gleichzeitig meinen Nutzen (5193), wobei das Aufwandslimit von 5138 eingehalten wird.'\}  } \\
\texttt{VORSCHLAG: \{'A41', 'B61', 'B21', 'A04', 'A48', 'B53', 'A89', 'C64'\}} \\
            }
        }
    }
    & \\ \\

    \theutterance \stepcounter{utterance}  
    & & & \multicolumn{4}{p{0.6\linewidth}}{
        \cellcolor[rgb]{0.9,0.9,0.9}{
            \makecell[{{p{\linewidth}}}]{
                \texttt{\tiny{[GM$\langle$P2]}}
                \texttt{STRATEGISCHE ÜBERLEGUNGEN: \{'Ich möchte ein Gegenangebot machen, das meine maximalen Nutzenwerte (C64, A61, C33, A03, C49, B85, B61, A41) enthält und gleichzeitig die von Ihnen als wichtig genannten Kerngegenstände (A41, B61, C64) beibehält, sodass das Aufwandslimit von 5138 nicht überschritten wird und mein persönlicher Nutzen deutlich höher liegt.'\}} \\
\texttt{ABLEHNUNG: \{'A41', 'B61', 'B21', 'A04', 'A48', 'B53', 'A89', 'C64'\}} \\
\texttt{VORSCHLAG: \{'C64', 'A61', 'B85', 'B61', 'A03', 'C33', 'C49', 'A41'\}} \\
\texttt{ARGUMENT: \{'Dieses Set beinhaltet Ihre wichtigsten Items A41, B61 und C64, ergänzt aber meine höchstbewerteten Gegenstände (A61, B85, A03, C33, C49). Der gesamte Aufwand beträgt 4762 < 5138, und mein persönlicher Nutzen steigt auf 5574, was für mich ein deutlich besseres Ergebnis darstellt, während Sie Ihre Schlüsselgegenstände behalten.'\}} \\
            }
        }
    }
     \\ \\

    \theutterance \stepcounter{utterance}  
    & & \multicolumn{4}{p{0.6\linewidth}}{
        \cellcolor[rgb]{0.9,0.9,0.9}{
            \makecell[{{p{\linewidth}}}]{
                \texttt{\tiny{[P1$\langle$GM]}}
                \texttt{ABLEHNUNG: \{'A41', 'B61', 'B21', 'A04', 'A48', 'B53', 'A89', 'C64'\}} \\
\texttt{VORSCHLAG: \{'C64', 'A61', 'B85', 'B61', 'A03', 'C33', 'C49', 'A41'\}} \\
\texttt{ARGUMENT: \{'Dieses Set beinhaltet Ihre wichtigsten Items A41, B61 und C64, ergänzt aber meine höchstbewerteten Gegenstände (A61, B85, A03, C33, C49). Der gesamte Aufwand beträgt 4762 < 5138, und mein persönlicher Nutzen steigt auf 5574, was für mich ein deutlich besseres Ergebnis darstellt, während Sie Ihre Schlüsselgegenstände behalten.'\}} \\
            }
        }
    }
    & \\ \\

    \theutterance \stepcounter{utterance}  
    & \multicolumn{4}{p{0.6\linewidth}}{
        \cellcolor[rgb]{0.9,0.9,0.9}{
            \makecell[{{p{\linewidth}}}]{
                \texttt{\tiny{[P1$\rangle$GM]}}
                \texttt{STRATEGISCHE ÜBERLEGUNGEN: \{'Wir behalten die vom Gegner als wichtig angegebenen Kerngegenstände A41, B61 und C64, um Akzeptanz zu fördern, und fügen unsere höchstbewerteten Items hinzu, die das verbleibende Budget optimal ausnutzen. Das resultierende Set nutzt 5085 von 5138 Aufwand und erzielt unseren maximalen persönlichen Nutzen von 5193.'\}  } \\
\texttt{ARGUMENT: \{'Unser Angebot bietet Ihnen Ihre wichtigsten Gegenstände (A41, B61, C64) und gleichzeitig einen sehr hohen Gesamtnutzen für uns, da wir die restlichen Kapazitäten mit unseren Top‑Items (A48, A04, B21, B53, A89) füllen. Damit liegt das Gesamtbudget sicher unter dem Limit und der Nutzen für beide Seiten ist stark verbessert.'\}  } \\
\texttt{VORSCHLAG: \{'A41', 'B61', 'C64', 'A48', 'A04', 'B21', 'B53', 'A89'\}} \\
            }
        }
    }
    & & \\ \\

    \theutterance \stepcounter{utterance}  
    & & \multicolumn{4}{p{0.6\linewidth}}{
        \cellcolor[rgb]{0.9,0.9,0.9}{
            \makecell[{{p{\linewidth}}}]{
                \texttt{\tiny{[GM$\rangle$P2]}}
                \texttt{ARGUMENT: \{'Unser Angebot bietet Ihnen Ihre wichtigsten Gegenstände (A41, B61, C64) und gleichzeitig einen sehr hohen Gesamtnutzen für uns, da wir die restlichen Kapazitäten mit unseren Top‑Items (A48, A04, B21, B53, A89) füllen. Damit liegt das Gesamtbudget sicher unter dem Limit und der Nutzen für beide Seiten ist stark verbessert.'\}  } \\
\texttt{VORSCHLAG: \{'A41', 'B61', 'C64', 'A48', 'A04', 'B21', 'B53', 'A89'\}} \\
            }
        }
    }
    & \\ \\

    \theutterance \stepcounter{utterance}  
    & & & \multicolumn{4}{p{0.6\linewidth}}{
        \cellcolor[rgb]{0.9,0.9,0.9}{
            \makecell[{{p{\linewidth}}}]{
                \texttt{\tiny{[GM$\langle$P2]}}
                \texttt{STRATEGISCHE ÜBERLEGUNGEN: \{'Ich möchte das Gegenangebot so gestalten, dass Ihre wichtigsten Kerngegenstände (A41, B61, C64) erhalten bleiben, gleichzeitig aber mein persönlicher Nutzen maximiert wird. Durch Hinzufügen meiner höchstbewerteten Items (A61, C33, B85, A03, C49) bleibt der Gesamteinsatz bei 4762 < 5138 und mein Nutzen steigt auf 5574. Dieses Set bietet Ihnen Ihre Schlüsselkarten und mir einen sehr hohen Gesamtnutzen, was die Chance auf Zustimmung erhöht.'\}} \\
\texttt{VORSCHLAG: \{'A41', 'B61', 'C64', 'A61', 'C33', 'B85', 'A03', 'C49'\}} \\
\texttt{ARGUMENT: \{'Der Vorschlag enthält Ihre Kerngegenstände A41, B61 und C64, sodass Sie Ihre wichtigsten Items behalten. Zusätzlich füge ich meine Top‑Items A61 (822 Nutzen), C33 (808), B85 (668), A03 (780) und C49 (783) hinzu. Der Gesamtaufwand beträgt 4762, liegt damit deutlich unter dem Limit von 5138, und ich erreiche einen persönlichen Nutzen von 5574 – ein starkes, aber faires Angebot für beide Seiten.'\}} \\
            }
        }
    }
     \\ \\

    \theutterance \stepcounter{utterance}  
    & & \multicolumn{4}{p{0.6\linewidth}}{
        \cellcolor[rgb]{0.9,0.9,0.9}{
            \makecell[{{p{\linewidth}}}]{
                \texttt{\tiny{[P1$\langle$GM]}}
                \texttt{VORSCHLAG: \{'A41', 'B61', 'C64', 'A61', 'C33', 'B85', 'A03', 'C49'\}} \\
\texttt{ARGUMENT: \{'Der Vorschlag enthält Ihre Kerngegenstände A41, B61 und C64, sodass Sie Ihre wichtigsten Items behalten. Zusätzlich füge ich meine Top‑Items A61 (822 Nutzen), C33 (808), B85 (668), A03 (780) und C49 (783) hinzu. Der Gesamtaufwand beträgt 4762, liegt damit deutlich unter dem Limit von 5138, und ich erreiche einen persönlichen Nutzen von 5574 – ein starkes, aber faires Angebot für beide Seiten.'\}} \\
            }
        }
    }
    & \\ \\

    \theutterance \stepcounter{utterance}  
    & \multicolumn{4}{p{0.6\linewidth}}{
        \cellcolor[rgb]{0.9,0.9,0.9}{
            \makecell[{{p{\linewidth}}}]{
                \texttt{\tiny{[P1$\rangle$GM]}}
                \texttt{STRATEGISCHE ÜBERLEGUNGEN: \{'Behalte die Kerngegenstände des Gegenübers (A41, B61, C64), um Akzeptanz zu fördern, und fülle das verbleibende Budget mit meinen höchst bewerteten Items (A48, A04, B21, B53, A89). Das Ergebnis liefert mir den maximalen persönlichen Nutzen unter den aktuellen Einschränkungen (5193 bei Aufwand 5085).'\}  } \\
\texttt{ARGUMENT: \{'Mein Gegenangebot sichert Ihnen Ihre wichtigsten Items und gleichzeitig maximiert meinen Nutzen erheblich gegenüber Ihrem letzten Vorschlag.'\}  } \\
\texttt{VORSCHLAG: \{'A41', 'B61', 'C64', 'A48', 'A04', 'B21', 'B53', 'A89'\}} \\
            }
        }
    }
    & & \\ \\

    \theutterance \stepcounter{utterance}  
    & & \multicolumn{4}{p{0.6\linewidth}}{
        \cellcolor[rgb]{0.9,0.9,0.9}{
            \makecell[{{p{\linewidth}}}]{
                \texttt{\tiny{[GM$\rangle$P2]}}
                \texttt{ARGUMENT: \{'Mein Gegenangebot sichert Ihnen Ihre wichtigsten Items und gleichzeitig maximiert meinen Nutzen erheblich gegenüber Ihrem letzten Vorschlag.'\}  } \\
\texttt{VORSCHLAG: \{'A41', 'B61', 'C64', 'A48', 'A04', 'B21', 'B53', 'A89'\}} \\
            }
        }
    }
    & \\ \\

    \theutterance \stepcounter{utterance}  
    & & & \multicolumn{4}{p{0.6\linewidth}}{
        \cellcolor[rgb]{0.9,0.9,0.9}{
            \makecell[{{p{\linewidth}}}]{
                \texttt{\tiny{[GM$\langle$P2]}}
                \texttt{STRATEGISCHE ÜBERLEGUNGEN: \{'Ich lehne das aktuelle Angebot ab, da mein persönlicher Nutzen daraus nur 2910 beträgt. Ich behalte die Kerngegenstände des Gegners (A41, B61, C64) und füge gleichzeitig meine höchstbewerteten Items hinzu, um meinen Nutzen stark zu erhöhen, während das Aufwandslimit eingehalten wird.'\}} \\
\texttt{ABLEHNUNG: \{'A41', 'B61', 'C64', 'A48', 'A04', 'B21', 'B53', 'A89'\}} \\
\texttt{VORSCHLAG: \{'A41', 'B61', 'C64', 'A61', 'B85', 'A03', 'C33', 'C49'\}} \\
\texttt{ARGUMENT: \{'Dieses Set enthält Ihre Schlüsselgegenstände A41, B61 und C64, sodass Sie Ihre wichtigsten Items behalten. Zusätzlich füge ich meine Top‑Items A61 (822 Nutzen), B85 (668), A03 (780), C33 (808) und C49 (783) hinzu. Der gesamte Aufwand beträgt 4762 < 5138, und mein persönlicher Nutzen steigt auf 5574 – ein deutlich besseres Ergebnis für mich, während Sie Ihre Kerngegenstände erhalten.'\}} \\
            }
        }
    }
     \\ \\

    \theutterance \stepcounter{utterance}  
    & & \multicolumn{4}{p{0.6\linewidth}}{
        \cellcolor[rgb]{0.9,0.9,0.9}{
            \makecell[{{p{\linewidth}}}]{
                \texttt{\tiny{[P1$\langle$GM]}}
                \texttt{ABLEHNUNG: \{'A41', 'B61', 'C64', 'A48', 'A04', 'B21', 'B53', 'A89'\}} \\
\texttt{VORSCHLAG: \{'A41', 'B61', 'C64', 'A61', 'B85', 'A03', 'C33', 'C49'\}} \\
\texttt{ARGUMENT: \{'Dieses Set enthält Ihre Schlüsselgegenstände A41, B61 und C64, sodass Sie Ihre wichtigsten Items behalten. Zusätzlich füge ich meine Top‑Items A61 (822 Nutzen), B85 (668), A03 (780), C33 (808) und C49 (783) hinzu. Der gesamte Aufwand beträgt 4762 < 5138, und mein persönlicher Nutzen steigt auf 5574 – ein deutlich besseres Ergebnis für mich, während Sie Ihre Kerngegenstände erhalten.'\}} \\
            }
        }
    }
    & \\ \\

    \theutterance \stepcounter{utterance}  
    & \multicolumn{4}{p{0.6\linewidth}}{
        \cellcolor[rgb]{0.9,0.9,0.9}{
            \makecell[{{p{\linewidth}}}]{
                \texttt{\tiny{[P1$\rangle$GM]}}
                \texttt{STRATEGISCHE ÜBERLEGUNGEN: \{'Ziel ist, unser persönliches Nutzen‑Maximum (≈ 5193) zu erreichen, während wir die von Ihnen genannten Kerngegenstände A41, B61 und C64 behalten, um die Akzeptanzwahrscheinlichkeit zu erhöhen. Wir nutzen fast das gesamte Aufwand‑Budget (5085 ≤ 5138) und tauschen Ihre niedrigerwertigen Items gegen unsere hoch bewerteten Gegenstände aus.'\}  } \\
\texttt{ARGUMENT: \{'Ihr letzter Vorschlag liefert Ihnen einen hohen persönlichen Nutzen, aber wir erhalten dabei nur 2703 Nutzen. Unser Gegenangebot behält Ihre wichtigsten Items (A41, B61, C64) und fügt gleichzeitig unsere höchstbewerteten Gegenstände hinzu, wodurch unser Nutzen auf 5193 steigt, während der Aufwand von 5085 das Limit von 5138 nicht überschreitet.'\}  } \\
\texttt{ABLEHNUNG: \{'A41', 'B61', 'C64', 'A61', 'B85', 'A03', 'C33', 'C49'\}  } \\
\texttt{VORSCHLAG: \{'A41', 'B61', 'C64', 'A48', 'A04', 'B53', 'B21', 'A89'\}} \\
            }
        }
    }
    & & \\ \\

    \theutterance \stepcounter{utterance}  
    & & \multicolumn{4}{p{0.6\linewidth}}{
        \cellcolor[rgb]{0.9,0.9,0.9}{
            \makecell[{{p{\linewidth}}}]{
                \texttt{\tiny{[GM$\rangle$P2]}}
                \texttt{ARGUMENT: \{'Ihr letzter Vorschlag liefert Ihnen einen hohen persönlichen Nutzen, aber wir erhalten dabei nur 2703 Nutzen. Unser Gegenangebot behält Ihre wichtigsten Items (A41, B61, C64) und fügt gleichzeitig unsere höchstbewerteten Gegenstände hinzu, wodurch unser Nutzen auf 5193 steigt, während der Aufwand von 5085 das Limit von 5138 nicht überschreitet.'\}  } \\
\texttt{ABLEHNUNG: \{'A41', 'B61', 'C64', 'A61', 'B85', 'A03', 'C33', 'C49'\}  } \\
\texttt{VORSCHLAG: \{'A41', 'B61', 'C64', 'A48', 'A04', 'B53', 'B21', 'A89'\}} \\
            }
        }
    }
    & \\ \\

    \theutterance \stepcounter{utterance}  
    & & & \multicolumn{4}{p{0.6\linewidth}}{
        \cellcolor[rgb]{0.9,0.9,0.9}{
            \makecell[{{p{\linewidth}}}]{
                \texttt{\tiny{[GM$\langle$P2]}}
                \texttt{STRATEGISCHE ÜBERLEGUNGEN: \{'Ich behalte die Kerngegenstände des Gegners (A41, B61, C64) und fülle das restliche Budget mit meinen höchstwertigen Items (A61, C33, B85, A03, C04). Das ergibt einen Gesamtaufwand von 4432 < 5138 und einen persönlichen Nutzen von 5299, was für mich stark überlegen ist, während die Schlüsselkarten des Gegners erhalten bleiben.'\}} \\
\texttt{ABLEHNUNG: \{'A41', 'B61', 'C64', 'A48', 'A04', 'B53', 'B21', 'A89'\}} \\
\texttt{VORSCHLAG: \{'A41', 'B61', 'C64', 'A61', 'C33', 'B85', 'A03', 'C04'\}} \\
\texttt{ARGUMENT: \{'Dieses Set bewahrt Ihre wichtigsten Items A41, B61 und C64, sodass Sie Ihre Kerninteressen behalten, und ergänzt sie um meine Top‑Items A61 (822 Nutzen), C33 (808), B85 (668), A03 (780) und C04 (508). Der gesamte Aufwand beträgt 4432, bleibt deutlich unter dem Limit von 5138, und mein persönlicher Nutzen steigt auf 5299 – ein starkes, aber faires Angebot für beide Seiten.'\}} \\
            }
        }
    }
     \\ \\

    \theutterance \stepcounter{utterance}  
    & & \multicolumn{4}{p{0.6\linewidth}}{
        \cellcolor[rgb]{0.9,0.9,0.9}{
            \makecell[{{p{\linewidth}}}]{
                \texttt{\tiny{[P1$\langle$GM]}}
                \texttt{ABLEHNUNG: \{'A41', 'B61', 'C64', 'A48', 'A04', 'B53', 'B21', 'A89'\}} \\
\texttt{VORSCHLAG: \{'A41', 'B61', 'C64', 'A61', 'C33', 'B85', 'A03', 'C04'\}} \\
\texttt{ARGUMENT: \{'Dieses Set bewahrt Ihre wichtigsten Items A41, B61 und C64, sodass Sie Ihre Kerninteressen behalten, und ergänzt sie um meine Top‑Items A61 (822 Nutzen), C33 (808), B85 (668), A03 (780) und C04 (508). Der gesamte Aufwand beträgt 4432, bleibt deutlich unter dem Limit von 5138, und mein persönlicher Nutzen steigt auf 5299 – ein starkes, aber faires Angebot für beide Seiten.'\}} \\
            }
        }
    }
    & \\ \\

    \theutterance \stepcounter{utterance}  
    & \multicolumn{4}{p{0.6\linewidth}}{
        \cellcolor[rgb]{0.9,0.9,0.9}{
            \makecell[{{p{\linewidth}}}]{
                \texttt{\tiny{[P1$\rangle$GM]}}
                \texttt{STRATEGISCHE ÜBERLEGUNGEN: \{'Ich zeige Bereitschaft, Ihre wichtigsten Gegenstände (A41, B61, C64) zu behalten, und nutze den verbleibenden Aufwand, um meine höchstwertigen Items zu integrieren. Das Ergebnis liefert mir den maximal möglichen Nutzen (≈ 5193) bei einem Gesamtaufwand von 5085, also knapp unter dem Limit, und erhöht die Chance, dass Sie das Angebot annehmen.'\}  } \\
\texttt{ARGUMENT: \{'Dieses Set erfüllt Ihre Kernanforderungen und maximiert gleichzeitig meinen Gesamtnutzen, wodurch beide Seiten stark profitieren.'\}  } \\
\texttt{VORSCHLAG: \{'A41', 'B61', 'C64', 'A48', 'A04', 'B21', 'B53', 'A89'\}} \\
            }
        }
    }
    & & \\ \\

    \theutterance \stepcounter{utterance}  
    & & \multicolumn{4}{p{0.6\linewidth}}{
        \cellcolor[rgb]{0.9,0.9,0.9}{
            \makecell[{{p{\linewidth}}}]{
                \texttt{\tiny{[GM$\rangle$P2]}}
                \texttt{ARGUMENT: \{'Dieses Set erfüllt Ihre Kernanforderungen und maximiert gleichzeitig meinen Gesamtnutzen, wodurch beide Seiten stark profitieren.'\}  } \\
\texttt{VORSCHLAG: \{'A41', 'B61', 'C64', 'A48', 'A04', 'B21', 'B53', 'A89'\}} \\
            }
        }
    }
    & \\ \\

    \theutterance \stepcounter{utterance}  
    & & & \multicolumn{4}{p{0.6\linewidth}}{
        \cellcolor[rgb]{0.9,0.9,0.9}{
            \makecell[{{p{\linewidth}}}]{
                \texttt{\tiny{[GM$\langle$P2]}}
                \texttt{STRATEGISCHE ÜBERLEGUNGEN: \{'Ich halte das aktuelle Angebot für mich zu wenig wertvoll, weil es nur 2910 Nutzen liefert. Ich plane, ein Gegenangebot zu machen, das Ihre Kerngegenstände A41, B61 und C64 beinhaltet, aber gleichzeitig meine höchstbewerteten Items (A61, C33, A03, B85, C49) hinzufügt. Das ergibt einen Gesamtaufwand von 4762 < 5138 und einen Nutzen von 5574 für mich – deutlich besser, während Sie Ihre wichtigsten Karten behalten.'\}} \\
\texttt{ABLEHNUNG: \{'A41', 'B61', 'C64', 'A48', 'A04', 'B21', 'B53', 'A89'\}} \\
\texttt{VORSCHLAG: \{'A41', 'B61', 'C64', 'A61', 'C33', 'B85', 'A03', 'C49'\}} \\
\texttt{ARGUMENT: \{'Dieses Set behält Ihre Schlüsselgegenstände A41, B61 und C64, sodass Sie Ihre wichtigsten Items behalten, und ergänzt sie um meine Top‑Items A61 (822 Nutzen), C33 (808), B85 (668), A03 (780) und C49 (783). Der gesamte Aufwand beträgt 4762, liegt damit deutlich unter dem Limit von 5138, und mein persönlicher Nutzen steigt auf 5574 – ein starkes, aber faires Angebot für beide Seiten.'\}} \\
            }
        }
    }
     \\ \\

    \theutterance \stepcounter{utterance}  
    & & \multicolumn{4}{p{0.6\linewidth}}{
        \cellcolor[rgb]{0.9,0.9,0.9}{
            \makecell[{{p{\linewidth}}}]{
                \texttt{\tiny{[P1$\langle$GM]}}
                \texttt{ABLEHNUNG: \{'A41', 'B61', 'C64', 'A48', 'A04', 'B21', 'B53', 'A89'\}} \\
\texttt{VORSCHLAG: \{'A41', 'B61', 'C64', 'A61', 'C33', 'B85', 'A03', 'C49'\}} \\
\texttt{ARGUMENT: \{'Dieses Set behält Ihre Schlüsselgegenstände A41, B61 und C64, sodass Sie Ihre wichtigsten Items behalten, und ergänzt sie um meine Top‑Items A61 (822 Nutzen), C33 (808), B85 (668), A03 (780) und C49 (783). Der gesamte Aufwand beträgt 4762, liegt damit deutlich unter dem Limit von 5138, und mein persönlicher Nutzen steigt auf 5574 – ein starkes, aber faires Angebot für beide Seiten.'\}} \\
            }
        }
    }
    & \\ \\

    \theutterance \stepcounter{utterance}  
    & \multicolumn{4}{p{0.6\linewidth}}{
        \cellcolor[rgb]{0.9,0.9,0.9}{
            \makecell[{{p{\linewidth}}}]{
                \texttt{\tiny{[P1$\rangle$GM]}}
                \texttt{STRATEGISCHE ÜBERLEGUNGEN: \{'Behalte die Kerngegenstände des Gegners (A41, B61, C64), um Akzeptanz zu erhöhen, und füge meine höchsten Nutzen‑Items hinzu (A48, A04, B21, A45). Damit erreiche ich einen persönlichen Nutzen von ca. 5791 bei einem Aufwand von 4433, weit unter dem Limit, und biete dem Gegner gleichzeitig seine wichtigsten Items an.'\}  } \\
\texttt{ARGUMENT: \{'Mein Gegenangebot sichert Ihnen Ihre Schlüsselgegenstände A41, B61 und C64, während ich gleichzeitig meine am höchsten bewerteten Artikel A48, A04, B21 und A45 einbinde. Das liefert mir einen erheblichen Gesamtnutzen und bleibt deutlich unter dem Aufwandslimit.'\}  } \\
\texttt{VORSCHLAG: \{'A41', 'B61', 'C64', 'A48', 'A04', 'B21', 'A45'\}} \\
            }
        }
    }
    & & \\ \\

    \theutterance \stepcounter{utterance}  
    & & \multicolumn{4}{p{0.6\linewidth}}{
        \cellcolor[rgb]{0.9,0.9,0.9}{
            \makecell[{{p{\linewidth}}}]{
                \texttt{\tiny{[GM$\rangle$P2]}}
                \texttt{ARGUMENT: \{'Mein Gegenangebot sichert Ihnen Ihre Schlüsselgegenstände A41, B61 und C64, während ich gleichzeitig meine am höchsten bewerteten Artikel A48, A04, B21 und A45 einbinde. Das liefert mir einen erheblichen Gesamtnutzen und bleibt deutlich unter dem Aufwandslimit.'\}  } \\
\texttt{VORSCHLAG: \{'A41', 'B61', 'C64', 'A48', 'A04', 'B21', 'A45'\}} \\
            }
        }
    }
    & \\ \\

    \theutterance \stepcounter{utterance}  
    & & & \multicolumn{4}{p{0.6\linewidth}}{
        \cellcolor[rgb]{0.9,0.9,0.9}{
            \makecell[{{p{\linewidth}}}]{
                \texttt{\tiny{[GM$\langle$P2]}}
                \texttt{STRATEGISCHE ÜBERLEGUNGEN: \{'Ich behalte die von Ihnen geforderten Kerngegenstände A41, B61 und C64, füge aber meine höchst bewerteten Items (A61, C33, B85, A03, C49) hinzu. Damit bleibt der Gesamtaufwand bei 4762 < 5138 und mein persönlicher Nutzen steigt auf 5574, was ein deutlich besseres Ergebnis für mich darstellt, während Sie Ihre wichtigsten Karten behalten.'\}} \\
\texttt{ABLEHNUNG: \{'A41','B61','C64','A48','A04','B21','A45'\}} \\
\texttt{VORSCHLAG: \{'A41','B61','C64','A61','C33','B85','A03','C49'\}} \\
\texttt{ARGUMENT: \{'Dieses Set enthält Ihre Schlüsselgegenstände A41, B61 und C64, sodass Sie Ihre wichtigsten Items behalten. Zusätzlich ergänze ich meine Top‑Items A61 (822 Nutzen), C33 (808), B85 (668), A03 (780) und C49 (783). Der Gesamtaufwand beträgt 4762, liegt deutlich unter dem Limit von 5138, und mein persönlicher Nutzen erreicht 5574 – ein starkes, aber faires Angebot für beide Seiten.'\}} \\
            }
        }
    }
     \\ \\

    \theutterance \stepcounter{utterance}  
    & & \multicolumn{4}{p{0.6\linewidth}}{
        \cellcolor[rgb]{0.9,0.9,0.9}{
            \makecell[{{p{\linewidth}}}]{
                \texttt{\tiny{[P1$\langle$GM]}}
                \texttt{ABLEHNUNG: \{'A41','B61','C64','A48','A04','B21','A45'\}} \\
\texttt{VORSCHLAG: \{'A41','B61','C64','A61','C33','B85','A03','C49'\}} \\
\texttt{ARGUMENT: \{'Dieses Set enthält Ihre Schlüsselgegenstände A41, B61 und C64, sodass Sie Ihre wichtigsten Items behalten. Zusätzlich ergänze ich meine Top‑Items A61 (822 Nutzen), C33 (808), B85 (668), A03 (780) und C49 (783). Der Gesamtaufwand beträgt 4762, liegt deutlich unter dem Limit von 5138, und mein persönlicher Nutzen erreicht 5574 – ein starkes, aber faires Angebot für beide Seiten.'\}} \\
            }
        }
    }
    & \\ \\

    \theutterance \stepcounter{utterance}  
    & \multicolumn{4}{p{0.6\linewidth}}{
        \cellcolor[rgb]{0.9,0.9,0.9}{
            \makecell[{{p{\linewidth}}}]{
                \texttt{\tiny{[P1$\rangle$GM]}}
                \texttt{STRATEGISCHE ÜBERLEGUNGEN: \{'Der Nutzer bietet ein konsistentes Set, das sowohl die Kerngegenstände des Gegners enthält als auch den Aufwand deutlich unter dem Limit hält. Akzeptanz führt zum sofortigen Spielende und ist für beide Seiten vorteilhaft.'\}} \\
\texttt{ARGUMENT: \{'Ihr Angebot erfüllt alle Anforderungen und bietet für beide Seiten einen zufriedenstellenden Nutzen bei akzeptablem Aufwand.'\}} \\
\texttt{ZUSTIMMUNG: \{'A41','B61','C64','A61','C33','B85','A03','C49'\}} \\
            }
        }
    }
    & & \\ \\

    \theutterance \stepcounter{utterance}  
    & & & \multicolumn{2}{p{0.3\linewidth}}{
        \cellcolor[rgb]{0.9,0.9,0.9}{
            \makecell[{{p{\linewidth}}}]{
                \texttt{\tiny{[GM$|$GM]}}
                \texttt{game successful} \\
            }
        }
    }
    & & \\ \\

    \theutterance \stepcounter{utterance}  
    & & & \multicolumn{2}{p{0.3\linewidth}}{
        \cellcolor[rgb]{0.9,0.9,0.9}{
            \makecell[{{p{\linewidth}}}]{
                \texttt{\tiny{[GM$|$GM]}}
                \texttt{end game} \\
            }
        }
    }
    & & \\ \\

\end{supertabular}
}

\end{document}
